\documentclass[10pt,a4paper]{amsart}
\usepackage[margin=19mm]{geometry}
\usepackage{amsmath,amssymb,mathtools}
\usepackage[scr=rsfs,cal=euler]{mathalfa}
\usepackage{xfrac}
\usepackage[unicode,hidelinks,bookmarks=false]{hyperref}
\newcommand{\ie}{i.e.\ }
\newcommand{\cf}{cf.\ }
\newcommand{\resp}{respectively\ }
\newcommand{\Subsection}{\subsection}
\providecommand{\nobreakdash}{\nobreak\hbox{-}}
\setlength{\emergencystretch}{2em}
\multlinegap0pt
\usepackage{bm}
\usepackage{bbm}
\usepackage{dsfont}
\usepackage{xspace}
\usepackage{cancel}
\usepackage{mathtools}
\usepackage{amsthm}
\makeatletter
\@ifundefined{lemma}{\newtheorem{lemma}{Lemma}}{}
\makeatother
\title[Extending Ghouila-Houri's Characterization of Comparability ]{Extending Ghouila-Houri's Characterization of Comparability Graphs to Temporal Graphs}
\author{Pierre Charbit, Michel Habib, Amalia Sorondo}
\date{Source: arXiv:2510.06849v3 (CC BY 4.0)}
\begin{document}
\maketitle

\noindent\emph{Source and licence.} Extending Ghouila-Houri's Characterization of Comparability Graphs to Temporal Graphs. Version arXiv:2510.06849v3 (CC BY 4.0), distributed under the Creative Commons Attribution 4.0 International licence, as recorded in the arXiv licence metadata for this exact version and shown on the abstract page https://arxiv.org/abs/2510.06849v3. Full rights notice: licenses/arxiv-2510.06849-CC-BY-4.0.txt.\par\smallskip

\noindent\textit{Editorial scope.} Counted unit: the Lemma whose source label is \texttt{None} in the article (source0.tex lines 789--791, source label \texttt{None}), delivered together with the article's own complete original proof of it (source0.tex lines 793--809, one \texttt{proof} environment of 17 lines). No mathematics is written, repaired or re-derived: statement and proof are the article's own text line for line. Required context retained uncounted: none. Every definition, display and label of those lines is retained unchanged. Omitted from this package and not redistributed: 1--788, 792--792, 810--1132. Nothing inside a retained excerpt is omitted or altered: every retained body is delivered exactly as the article's lines are written, so no omission receipt of any kind accompanies this package. Citations: the retained excerpts carry no citation command at all, so no bibliography entry of the article is transcribed and no reference is redistributed. Reference adaptation: none was needed and none was made; every reference inside the delivered unit resolves inside it, so no unresolved reference or citation is produced. Printed slips kept literal: no printed word, symbol, hypothesis or number of the counted statement or of its proof was altered. Layout and class: the article's own class is \texttt{llncs} with packages including geometry, inputenc, fontenc, graphicx, amssymb, amsmath, dsfont, enumitem, natbib, tikz, hyperref, cleveref, xcolor, upgreek; the wrapper is the lane's pinned \texttt{amsart} editorial wrapper at 10pt on A4, which restates the theorem-like environments the retained text needs and adds the article's own macro definitions that the retained text needs (none), with extra wrapper packages (bm, bbm, dsfont, xspace, cancel, mathtools). The delivered numbering is the unit's own. Assets: no retained span contains a figure environment, a graphics-inclusion command, a PDF-page inclusion, a table or a TikZ picture; no image file is copied into this package, the package declares no asset dependency and no image file is redistributed with it. Licence: version arXiv:2510.06849v3 of this article is distributed under the Creative Commons Attribution 4.0 International licence, as recorded in the arXiv licence metadata for this exact version and shown on the abstract page https://arxiv.org/abs/2510.06849v3. A full rights notice is delivered at licenses/arxiv-2510.06849-CC-BY-4.0.txt. Characters: every retained excerpt is pure ASCII, so the delivered engine, XeLaTeX with its default Latin Modern text font, reports no missing character.
\begin{lemma}
If there exists a quasi-transitive orientation of $G$, then there also exists a transitive orientation of $G$.
\end{lemma}

\begin{proof}
We assume the lemma holds for $|X| < n$, and consider a graph $G = (X,U)$ with $|X| = n$ that admits a quasi-transitive orientation $V$. Observe that if three vertices $a,b,c$ satisfy $(a,b) \in V, (b,c) \in V, (c,a) \in V$, then any other vertex of $G$ is adjacent to 0, 2, or 3 of these vertices.

If $V$ is not transitive, then there exist three vertices $x_1, x_2, x_3$ such that $(x_1,x_2) \in V,(x_2,x_3) \in V,(x_3,x_1) \in V$. Two cases arise.

\paragraph{Case 1.}
Every vertex different from $x_1,x_2,x_3$ that is adjacent to one of them is adjacent to all three.

Remove $x_2$ and $x_3$, and orient the remaining graph transitively (this is possible by induction). Then orient each edge $\{x,x_2\}$ or $\{x,x_3\}$ (for $x \notin \{x_1,x_2,x_3\}$) in the same direction as $\{x,x_1\}$, and orient the triangle formed by $x_1,x_2,x_3$ transitively. This yields a transitive orientation of $G$.

\paragraph{Case 2.}
There exists at least one vertex different from $x_1,x_2,x_3$ adjacent to two of them but not the third; for example, adjacent to $x_2$ and $x_3$, but not to $x_1$. Let $A$ be the set of vertices different from $x_2$ and $x_3$ adjacent to both $x_2$ and $x_3$, let $W$ be the set of pairs of vertices in $A$ that are not edges of $G$. Let $C$ be the connected component containing $x_1$ in the graph $(A,W)$; then $|C| > 1$.

Every $c \in C$ satisfies $(c,x_2) \in V $ and $(x_3,c) \in V$. We show that any vertex not in $C$ that is adjacent to one vertex of $C$ is necessarily adjacent to all vertices of $C$. Indeed, let $x \notin C$ and $a,b \in C$. If $x$ is adjacent to $a$, then it must be adjacent to $x_2$ or $x_3$, hence $x$ also adjacent to $b$; otherwise we get $x \in A$ and $\{x,b\} \in W$, implying $x \in C$, a contradiction.

It is now sufficient to orient transitively the subgraph induced by $C$ first, then the graph obtained by removing all vertices of $C$ except $x_1$ and to orient every edge $\{x,x'\}$ with $x \notin C$ and $x' \in C$ in the same direction as $\{x,x_1\}$, to obtain a transitive orientation of $G$.
$\qed$\end{proof}

\end{document}
